Deux longs fils rectilignes et parallèles sont séparés d'une distance $d = \qty{\v{d}}{cm}$. Ils sont parcourus dans le même sens par des courants $I_1 = \qty{\v{I1}}{A}$ et $I_2 = \qty{\v{I2}}{A}$. ($\mu_0 = 4\pi \times \qty{e-7}{T.m/A}$)

\begin{enumerate}[label=\alph*)]
    \item Déterminez la grandeur du champ magnétique au point situé à mi-chemin entre les deux fils.
    \item Déterminez la force par unité de longueur qu'exerce un fil sur l'autre. Est-elle attractive ou répulsive ?
    \item À quelle distance du fil 1, entre les deux fils, le champ magnétique est-il nul ?
\end{enumerate}
